\section{Related Work}
\label{sec:related_work}

Our work builds on parameter merging, adaptive ensembling, and statistical learning theory.

\textbf{Weight-Space Model Merging}. Weight-space merging combines multiple fine-tuned models derived from a common pre-trained base model without training them from scratch. Linear parameter interpolation and Task Arithmetic~\cite{wortsman2022model} add task-specific parameter updates (task vectors) to the base model. To mitigate parameter interference, TIES-Merging~\cite{yadav23ties} prunes redundant parameters and resolves sign conflicts, while DARE-Merging~\cite{yu23dare} uses random dropout to prune task vectors up to $99\%$ of the weights. Other methods like Git Re-Basin~\cite{ainsworth2022git} and ZipIt!~\cite{gallego2024mergekit} align weight permutations before merging. These methods are static, setting uniform weights across all layers.

\textbf{Adaptive and Layer-Wise Ensembling}. Recent work shows that task vector influence varies significantly across layers~\cite{choshen22fusing}. Layer-wise ensembling dynamically tunes merging coefficients $\alpha_k(l)$ for each layer. For example, AdaMerging~\cite{yang2024adamerging} optimizes layer-wise weights using prediction entropy minimization during test-time adaptation. PolyMerge~\cite{yang2024adamerging} fits a continuous curve to the coefficients across layers. While these methods improve multi-task performance, they are highly overparameterized. When tuned on small calibration datasets (few-shot adaptation), unconstrained optimization overfits local sample noise, leading to transductive overfitting and poor out-of-distribution (OOD) generalization.

\textbf{Theoretical Foundations and Trajectory Bounding}. To address the overfitting of layer-wise tuning, researchers have introduced statistical learning theory. PAC-ZCA~\cite{huang2026pac} bounds temperature-only routing policies using PAC-Bayesian bounds, while PAC-Kinetics~\cite{gomez2026pac} models parameter dynamics as a chemical kinetics system. Closest to our work, Rademacher-Bounded Polynomial Merging (RBPM) restricts layer-wise coefficients to a low-degree polynomial subspace and applies an analytical Rademacher penalty to control hypothesis class capacity~\cite{chatterjee2024robustness}. However, polynomial trajectories of degree $d \ge 2$ suffer from Runge's-like boundary runaway, causing large oscillations and extreme values at the first and last layers which harms performance. Our proposed \textbf{RB-FTM} and \textbf{RB-DCTM} address this issue by projecting ensembling trajectories onto bounded continuous sinusoidal and cosinusoidal bases, ensuring boundary stability while providing tight empirical Rademacher complexity guarantees.